\documentclass[10pt,a4paper]{amsart}
\usepackage[margin=19mm]{geometry}
\usepackage{amsmath,amssymb,mathtools}
\usepackage[scr=rsfs,cal=euler]{mathalfa}
\usepackage{xfrac}
\usepackage[unicode,hidelinks,bookmarks=false]{hyperref}
\newcommand{\ie}{i.e.\ }
\newcommand{\cf}{cf.\ }
\newcommand{\resp}{respectively\ }
\newcommand{\Subsection}{\subsection}
\providecommand{\nobreakdash}{\nobreak\hbox{-}}
\setlength{\emergencystretch}{2em}
\multlinegap0pt
\usepackage{bm}
\usepackage{bbm}
\usepackage{dsfont}
\usepackage{xspace}
\usepackage{cancel}
\usepackage{mathtools}
\usepackage{amsthm}
\makeatletter
\@ifundefined{thm}{\newtheorem{thm}{Thm}}{}
\@ifundefined{prop}{\newtheorem{prop}[thm]{Proposition}}{}
\makeatother
\DeclareMathOperator{\diag}{diag}
\newcommand\set[1]{\left\{#1\right\}}
\title[Smith Normal Forms of Graphical Hermite Simplices]{Smith Normal Forms of Graphical Hermite Simplices}
\author{Benjamin Braun, Antwon Park}
\date{Source: arXiv:2511.03822v2 (CC BY 4.0)}
\begin{document}
\maketitle

\noindent\emph{Source and licence.} Smith Normal Forms of Graphical Hermite Simplices. Version arXiv:2511.03822v2 (CC BY 4.0), distributed under the Creative Commons Attribution 4.0 International licence, as recorded in the arXiv licence metadata for this exact version and shown on the abstract page https://arxiv.org/abs/2511.03822v2. Full rights notice: licenses/arxiv-2511.03822-CC-BY-4.0.txt.\par\smallskip

\noindent\textit{Editorial scope.} Counted unit: the Prop whose source label is \texttt{thm:bipartite\_elementary\_divisors} in the article (source0.tex lines 856--865, source label \texttt{thm:bipartite\_elementary\_divisors}), delivered together with the article's own complete original proof of it (source0.tex lines 867--913, one \texttt{proof} environment of 47 lines). No mathematics is written, repaired or re-derived: statement and proof are the article's own text line for line. Required context retained uncounted: none. Every definition, display and label of those lines is retained unchanged. Omitted from this package and not redistributed: 1--855, 866--866, 914--952. Nothing inside a retained excerpt is omitted or altered: every retained body is delivered exactly as the article's lines are written, so no omission receipt of any kind accompanies this package. Citations: the retained excerpts carry no citation command at all, so no bibliography entry of the article is transcribed and no reference is redistributed. Reference adaptation: none was needed and none was made; every reference inside the delivered unit resolves inside it, so no unresolved reference or citation is produced. Printed slips kept literal: no printed word, symbol, hypothesis or number of the counted statement or of its proof was altered. Layout and class: the article's own class is \texttt{amsart} with packages including geometry, amssymb, amsmath, graphicx, tikz, xcolor, algpseudocode, algorithm, todonotes; the wrapper is the lane's pinned \texttt{amsart} editorial wrapper at 10pt on A4, which restates the theorem-like environments the retained text needs and adds the article's own macro definitions that the retained text needs (\texttt{\textbackslash diag}, \texttt{\textbackslash set}), with extra wrapper packages (bm, bbm, dsfont, xspace, cancel, mathtools). The delivered numbering is the unit's own. Assets: no retained span contains a figure environment, a graphics-inclusion command, a PDF-page inclusion, a table or a TikZ picture; no image file is copied into this package, the package declares no asset dependency and no image file is redistributed with it. Licence: version arXiv:2511.03822v2 of this article is distributed under the Creative Commons Attribution 4.0 International licence, as recorded in the arXiv licence metadata for this exact version and shown on the abstract page https://arxiv.org/abs/2511.03822v2. A full rights notice is delivered at licenses/arxiv-2511.03822-CC-BY-4.0.txt. Characters: every retained excerpt is pure ASCII, so the delivered engine, XeLaTeX with its default Latin Modern text font, reports no missing character.
\begin{prop}\label{thm:bipartite_elementary_divisors}
	Let G be a bipartite DAG and B be given as 
	\[
	B = \sum_{(i,j)\in E(G)} E_{i,j}. 
	\]
	
	Let $\beta_i$ be the elementary divisors of $B$ and let $\gamma_i := \gcd(m, \beta_i)$. 
	If the rank of $B$ is $r$, then the SNF of $A_{\vec{m}, G}$ is
	\[\diag\left(\gamma_1, \gamma_2, \dots, \gamma_r, m, m, \dots, m, \dfrac{m^2}{\gamma_r}, \dfrac{m^2}{\gamma_{r-1}}, \dots, \dfrac{m^2}{\gamma_{1}}\right).\]
\end{prop}

\begin{proof}
	Without loss of generality, assume that $V(G) = \set{1,2,\dots,k} \amalg \set{k+1,\dots,n}$. 
	Then we have $x_{i,j}$ is nonzero only if $i \leq k$ and $j > k$. Thus, we have that $A_{\vec{m}, G}$ is a block matrix of the form
	\[A_{\vec{m}, G} = \begin{bmatrix}
		mI_k & B \\
		0 & mI_{n-k}
	\end{bmatrix}.\]
	
	The following discussion demonstrates that elementary row operations on $B$ can arise from elementary row operations on $A_{\vec{m} ,G}$ without affecting the other entries in $A_{\vec{m} ,G}$. 
	An identical argument will show the same applies for elementary column operations. 
	Let $i,j \leq k$ be distinct.
	\begin{itemize}
		\item To perform the row operation $R_i \to R_i + R_j$ in $B$, we can perform this operation in $A_{\vec{m}, G}$. In the $mI_k$ block, we are changing the $[i,j]$ position from 0 to $m$. 
		Note that column $i$ of $A_{\vec{m}, G}$ has an $m$ in the $i$-th row and 0 elsewhere. 
		We can then perform the column operation $C_j \to C_j - C_i$ such that $A_{\vec{m}, G}[i,j] = 0$. 
		\item To negate the $i$-th row in $B$, we can negate the $i$-th row in $A_{\vec{m}, G}$ and then negate the $i$-th column in $A_{\vec{m}, G}$.
		\item To swap the $i$-th and $j$-th rows in $B$, we can swap the $i$-th and $j$-th rows in $A_{\vec{m}, G}$. We can then swap the $i$-th and $j$-th columns so that the top left block is unaltered.
	\end{itemize}
	
	Therefore, we know it is possible to arrive at the SNF of $B$ via elementary row and column operations on $B$, so let us perform the corresponding sequence of elementary row and columns operations on $A_{\vec{m}, G}$ until we arrive at 
	\[A_{\vec{p}, G} \to \dots \to \begin{bmatrix}
		mI_k & \diag(\beta_1,\beta_2, \dots, \beta_r, 0, \dots, 0) \\
		0 & mI_{n-k}.
	\end{bmatrix}\]
	
	We can then perform a sequence of row and column permutations of this matrix such that the resulting matrix has $r$ $2\times 2$ blocks along the main diagonal of the form
	\[
	\begin{bmatrix}
		m & \beta_i \\ 0 & m
	\end{bmatrix},
	\]
	followed by a constant $m$ diagonal. 
	Specifically, exchanging columns $2$ and $k+1$ and exchanging rows $2$ and $k+1$ will place a $2\times 2$ block with $\beta_1$ in the upper left corner of $mI_k$ and will place a $2\times 2$ block with $\beta_2$ in the upper left corner of $mI_{n-k}$.
	Iterating this process for consecutive pairs of $\beta$'s will result in the desired form. 	
	
	For each positive $i \leq r$, there exists a sequence of elementary row and column operations (essentially, applying the Euclidean algorithm) such that
	\[\begin{bmatrix}
		m & \beta_i \\ 0 & m
	\end{bmatrix} \to \dots \to  \begin{bmatrix}
		\gamma_i& 0 \\ 0 & \dfrac{m^2}{\gamma_i}
	\end{bmatrix}.\]
	
	Therefore, $A_{\vec{m}, G}$ has the same SNF as 
	\[\diag\left(\gamma_1, \gamma_2, \dots, \gamma_r, m, m, \dots, m, \dfrac{m^2}{\gamma_r}, \dfrac{m^2}{\gamma_{r-1}}, \dots, \dfrac{m^2}{\gamma_{1}}\right).\]
	
	It is straightforward to verify from the divisibility property of the $\beta$'s that each of these terms divides the subsequent term, hence this is the SNF of $A_{\vec{m}, G}$.	
\end{proof}

\end{document}
